\begin{boxedabstract}
\small
Molecular dynamics (MD) simulations of metallic alloys need an interatomic
potential that describes every constituent element pair. The Modified
Embedded Atom Method (MEAM) provides one, but published MEAM parameter sets
each cover only a narrow slice of the periodic table, so an alloy whose
elements are split across several files cannot be simulated at all. This
project builds an automated pipeline that produces a single MEAM parameter
file for a chosen list of elements. The five stages are: (i) randomized
alloy configurations drawn from the available composition files, (ii) a
LAMMPS strain sweep that measures the cubic elastic constants and the derived
Young's modulus $E$ and Poisson's ratio $\nu$, (iii) density functional
theory (DFT) reference data computed in Quantum Espresso through the Atomic
Simulation Environment, (iv) DFT-driven MEAM initialization that merges the
existing libraries and fills the missing cross-terms, and (v) a neural-network
surrogate of the MEAM-to-$(C_{11},C_{12})$ map that is fed into the
Adam optimizer to refit the cross-interaction parameters. A separate composition-only
surrogate that maps element fractions straight to $(C_{11},C_{12})$ runs
alongside Stage~2 as a fast consistency check on the LAMMPS dataset and as a predictor for new compositions. On
42 named alloys with experimental data from ASM and MatWeb, both surrogates
reproduce the aluminium Young's modulus to within a few percent (composition
NN $2.3\,\%$, refitted MEAM $7.4\,\%$); accuracy degrades once the
composition leaves the aluminium-rich region the training data covers most
densely. The closing section sets out the sampling, loss, and
domain-coverage work that would extend the method past that region.
\vspace{0.5em}
\end{boxedabstract}
